\subsection{A current claim about Jones detection and its audited prerequisites}
\label{sec:jones-status}

The literature check was made on 8 October 2026. Carmi and Cohen's
June 2026 preprint explicitly claims that the Jones polynomial detects the
unknot \cite[Theorem~6.9]{carmicohen2026}. The arXiv record inspected on that
date lists version~1 only and displays no withdrawal notice. This claim
must be acknowledged; it cannot be resolved by repeating older descriptions
of the problem's status. Our theorems and recognition decisions do not assume
the claimed detection theorem.

There is a concrete reason for this independence. The pinned repository
already contains an audit of a prerequisite consequence of that preprint.
Lemma~5.8, combined with the universal realization asserted in
Proposition~5.12 and the threshold in Section~5.1.1, implies a uniform bound
on the degree in $p=x-x^{-1}$ of every ordinary residual twist state:
\[
 \deg_p V\le B_{\rm host}+11+2.
\]
Here the expansion is $V(x)=a(p)+b(p)x$, and the degree means
$\max(\deg a,\deg b)$. This implication precedes the low-order anchor
vanishing used later in Theorem~5.10; it is not restricted to hosts already
known to have trivial Jones polynomial \cite[Sections~5.1--5.6]{carmicohen2026}.

For clarity, the existing witness can be described without any clasp-machine
isotopy. Take the three-crossing host with planar-diagram rows
\[
 (5,1,2,0),\quad(1,4,3,2),\quad(0,3,4,5).
\]
Replace the third row by an odd anti-parallel twist of length $L$, keeping
its exterior ports fixed. Explicitly, the first twist row begins with
$(a_0,b_0)=(0,3)$. At row $i=0,\ldots,L-1$ use
$(a_i,b_i,c_i,d_i)$, where $(c_i,d_i)=(6+2i,7+2i)$ for $i<L-1$,
$(c_{L-1},d_{L-1})=(4,5)$, and
$(a_{i+1},b_{i+1})=(d_i,c_i)$. These are the same fixed-exterior
diagrams used by the baseline audit.

Their knot Jones polynomials, in $t=x^2$, satisfy the audited skein recurrence
\[
 V_1(t)=t+t^3-t^4,\qquad
 V_L(t)=t^2V_{L-2}(t)+t-t^2+t^3-t^4.
\]
Consequently the highest term is $-t^{L+3}$. The identity
$x^m=F_{m-1}(p)+F_m(p)x$, with
\[
 F_m(p)=\sum_{k=0}^{\lfloor(m-1)/2\rfloor}
          \binom{m-k-1}{k}p^{m-2k-1},
\]
shows that the corresponding $b$-polynomial has leading term
$-p^{2L+5}$. All lower $t$-powers contribute smaller $p$-degrees.
Thus the degree is $2L+5$; for the mirror it is $2L+6$.

The host has $B_{\rm host}=3\cdot3+1-1=9$, so the proposed uniform
residual cap is $22$. The baseline's $L=9$ witness has degree $23$.
The present package independently recomputes its complete Jones polynomial
with Regina~7.4.1, and recomputes the formal expansion using the displayed
binomial formula rather than the baseline's coefficient recurrence.
It also checks $L=23$: the degree is $51$, while $L>22$ even satisfies
the long-twist inequality needed by the later descent. The mirror degrees
are $24$ and $52$, respectively. All four polynomial computations agree
coefficient by coefficient with all three tensor arithmetic modes.

The original audit additionally records local oriented-smoothing reductions
and Reidemeister-II anchor cancellations. Our new independent computation
replicates its polynomial obstruction. It neither constructs a nontrivial
knot with Jones polynomial one nor identifies a unique erroneous
clasp-machine step. It shows that the stated universal geometric realization
and uniform residual-degree consequence cannot both apply to these explicit
ordinary twists. Accordingly, a Jones polynomial different from one gives
the usual nontriviality obstruction, whereas identity remains inconclusive
in this implementation. The source and complete raw outputs are
\path{reproduce/audit_regina_jones.py} and
\path{results/regina_jones_20261008.json}.
